\begin{figure}[htbp]
\centering
\begin{adjustbox}{width=\textwidth}
\begin{tikzpicture}[
    node distance=1.15cm,
    every node/.style={font=\small},
    layer/.style={
        draw,
        rounded corners,
        align=center,
        minimum width=2.0cm,
        minimum height=0.9cm
    },
    reg/.style={
        draw,
        dashed,
        rounded corners,
        align=center,
        minimum width=1.8cm,
        minimum height=0.75cm,
        font=\footnotesize
    },
    arrow/.style={-{Latex[length=2mm]}, thick}
]

\node[layer] (input) {Input\\$N$ features};
\node[reg, right=of input] (noise) {Gaussian\\noise};
\node[layer, right=of noise] (enc1) {Dense 36\\ReLU};
\node[reg, right=of enc1] (drop1) {Dropout};
\node[layer, right=of drop1] (bottleneck) {Dense 6\\ReLU\\Bottleneck};
\node[reg, right=of bottleneck] (drop2) {Dropout};
\node[layer, right=of drop2] (dec1) {Dense 36\\ReLU};
\node[layer, right=of dec1] (output) {Output\\$N$ features\\tanh};

\draw[arrow] (input) -- (noise);
\draw[arrow] (noise) -- (enc1);
\draw[arrow] (enc1) -- (drop1);
\draw[arrow] (drop1) -- (bottleneck);
\draw[arrow] (bottleneck) -- (drop2);
\draw[arrow] (drop2) -- (dec1);
\draw[arrow] (dec1) -- (output);

\node[above=0.25cm of enc1] {\textbf{Encoder}};
\node[above=0.25cm of dec1] {\textbf{Decoder}};

\end{tikzpicture}
\end{adjustbox}
\caption{CPS-GUARD-style autoencoder architecture used for reconstruction-based anomaly detection. The \(N\)-dimensional feature vector is regularized using Gaussian noise, compressed into a lower-dimensional bottleneck representation, and then reconstructed at the output. Dropout is applied after the first two hidden layers during training.}
\label{fig:cpsguard_autoencoder_architecture}
\end{figure}